% SHOP decomposition of the generic HTN -- include with % SHOP decomposition of the generic HTN -- include with % SHOP decomposition of the generic HTN -- include with % SHOP decomposition of the generic HTN -- include with \input{chapters/shop_example}
% Uses the compound/primitive/method styles and \htnpre from main.tex
\begin{figure}[htbp]
  \centering
  \tikzset{
    order/.style  ={circle, fill=green!45!black, text=white, font=\scriptsize\bfseries,
                    inner sep=0pt, minimum size=12pt},
    notused/.style={draw=black!30, text=black!40, fill=black!3},
  }
  \begin{forest}
    for tree={
      font=\small, align=center, inner sep=3pt, rounded corners=2pt,
      minimum width=1.0cm, minimum height=0.7cm,
      l sep=0.9cm, s sep=0.5cm,
      edge={draw=green!45!black, very thick},
      parent anchor=south, child anchor=north,
    }
    [{$T_1$}, compound
      [{$m_1$\htnpre{$p_1$}}, method
        [{$t_1$}, primitive, label={[order]below:1}]
        [{$T_2$}, compound
          [{$m_2$\htnpre{$p_2$}}, method, notused, edge={draw=black!30, thick, dashed}
            [{$t_2$}, primitive, notused, edge={draw=black!30, thick, dashed}]
          ]
          [{$m_3$\htnpre{$\neg p_2$}}, method
            [{$t_3$}, primitive, label={[order]below:2}]
            [{$t_4$}, primitive, label={[order]below:3}]
          ]
        ]
        [{$t_5$}, primitive, label={[order]below:4}]
      ]
    ]
  \end{forest}

  \medskip
  {\small Initial state $s_0 = \{p_1\}$ \quad$\Rightarrow$\quad
   plan $\pi = \langle t_1, t_3, t_4, t_5 \rangle$}

  \caption[SHOP decomposition of the generic HTN]{%
    Decomposition of the generic HTN by SHOP for the initial state
    $s_0 = \{p_1\}$. Highlighted edges show the chosen decomposition;
    method $m_2$ is not applicable because $p_2$ does not hold. Numbers
    give the order of the primitive tasks in the resulting plan.}
  \label{fig:shop-example}
\end{figure}

% Uses the compound/primitive/method styles and \htnpre from main.tex
\begin{figure}[htbp]
  \centering
  \tikzset{
    order/.style  ={circle, fill=green!45!black, text=white, font=\scriptsize\bfseries,
                    inner sep=0pt, minimum size=12pt},
    notused/.style={draw=black!30, text=black!40, fill=black!3},
  }
  \begin{forest}
    for tree={
      font=\small, align=center, inner sep=3pt, rounded corners=2pt,
      minimum width=1.0cm, minimum height=0.7cm,
      l sep=0.9cm, s sep=0.5cm,
      edge={draw=green!45!black, very thick},
      parent anchor=south, child anchor=north,
    }
    [{$T_1$}, compound
      [{$m_1$\htnpre{$p_1$}}, method
        [{$t_1$}, primitive, label={[order]below:1}]
        [{$T_2$}, compound
          [{$m_2$\htnpre{$p_2$}}, method, notused, edge={draw=black!30, thick, dashed}
            [{$t_2$}, primitive, notused, edge={draw=black!30, thick, dashed}]
          ]
          [{$m_3$\htnpre{$\neg p_2$}}, method
            [{$t_3$}, primitive, label={[order]below:2}]
            [{$t_4$}, primitive, label={[order]below:3}]
          ]
        ]
        [{$t_5$}, primitive, label={[order]below:4}]
      ]
    ]
  \end{forest}

  \medskip
  {\small Initial state $s_0 = \{p_1\}$ \quad$\Rightarrow$\quad
   plan $\pi = \langle t_1, t_3, t_4, t_5 \rangle$}

  \caption[SHOP decomposition of the generic HTN]{%
    Decomposition of the generic HTN by SHOP for the initial state
    $s_0 = \{p_1\}$. Highlighted edges show the chosen decomposition;
    method $m_2$ is not applicable because $p_2$ does not hold. Numbers
    give the order of the primitive tasks in the resulting plan.}
  \label{fig:shop-example}
\end{figure}

% Uses the compound/primitive/method styles and \htnpre from main.tex
\begin{figure}[htbp]
  \centering
  \tikzset{
    order/.style  ={circle, fill=green!45!black, text=white, font=\scriptsize\bfseries,
                    inner sep=0pt, minimum size=12pt},
    notused/.style={draw=black!30, text=black!40, fill=black!3},
  }
  \begin{forest}
    for tree={
      font=\small, align=center, inner sep=3pt, rounded corners=2pt,
      minimum width=1.0cm, minimum height=0.7cm,
      l sep=0.9cm, s sep=0.5cm,
      edge={draw=green!45!black, very thick},
      parent anchor=south, child anchor=north,
    }
    [{$T_1$}, compound
      [{$m_1$\htnpre{$p_1$}}, method
        [{$t_1$}, primitive, label={[order]below:1}]
        [{$T_2$}, compound
          [{$m_2$\htnpre{$p_2$}}, method, notused, edge={draw=black!30, thick, dashed}
            [{$t_2$}, primitive, notused, edge={draw=black!30, thick, dashed}]
          ]
          [{$m_3$\htnpre{$\neg p_2$}}, method
            [{$t_3$}, primitive, label={[order]below:2}]
            [{$t_4$}, primitive, label={[order]below:3}]
          ]
        ]
        [{$t_5$}, primitive, label={[order]below:4}]
      ]
    ]
  \end{forest}

  \medskip
  {\small Initial state $s_0 = \{p_1\}$ \quad$\Rightarrow$\quad
   plan $\pi = \langle t_1, t_3, t_4, t_5 \rangle$}

  \caption[SHOP decomposition of the generic HTN]{%
    Decomposition of the generic HTN by SHOP for the initial state
    $s_0 = \{p_1\}$. Highlighted edges show the chosen decomposition;
    method $m_2$ is not applicable because $p_2$ does not hold. Numbers
    give the order of the primitive tasks in the resulting plan.}
  \label{fig:shop-example}
\end{figure}

% Uses the compound/primitive/method styles and \htnpre from main.tex
\begin{figure}[htbp]
  \centering
  \tikzset{
    order/.style  ={circle, fill=green!45!black, text=white, font=\scriptsize\bfseries,
                    inner sep=0pt, minimum size=12pt},
    notused/.style={draw=black!30, text=black!40, fill=black!3},
  }
  \begin{forest}
    for tree={
      font=\small, align=center, inner sep=3pt, rounded corners=2pt,
      minimum width=1.0cm, minimum height=0.7cm,
      l sep=0.9cm, s sep=0.5cm,
      edge={draw=green!45!black, very thick},
      parent anchor=south, child anchor=north,
    }
    [{$T_1$}, compound
      [{$m_1$\htnpre{$p_1$}}, method
        [{$t_1$}, primitive, label={[order]below:1}]
        [{$T_2$}, compound
          [{$m_2$\htnpre{$p_2$}}, method, notused, edge={draw=black!30, thick, dashed}
            [{$t_2$}, primitive, notused, edge={draw=black!30, thick, dashed}]
          ]
          [{$m_3$\htnpre{$\neg p_2$}}, method
            [{$t_3$}, primitive, label={[order]below:2}]
            [{$t_4$}, primitive, label={[order]below:3}]
          ]
        ]
        [{$t_5$}, primitive, label={[order]below:4}]
      ]
    ]
  \end{forest}

  \medskip
  {\small Initial state $s_0 = \{p_1\}$ \quad$\Rightarrow$\quad
   plan $\pi = \langle t_1, t_3, t_4, t_5 \rangle$}

  \caption[SHOP decomposition of the generic HTN]{%
    Decomposition of the generic HTN by SHOP for the initial state
    $s_0 = \{p_1\}$. Highlighted edges show the chosen decomposition;
    method $m_2$ is not applicable because $p_2$ does not hold. Numbers
    give the order of the primitive tasks in the resulting plan.}
  \label{fig:shop-example}
\end{figure}
